\documentclass[11pt,a4paper]{article}
\usepackage[utf8]{inputenc}
\usepackage[T1]{fontenc}
\usepackage[brazil]{babel}
\usepackage[margin=2.5cm]{geometry}
\usepackage{mathptmx}
\usepackage{enumitem}
\usepackage{booktabs}
\usepackage{tabularx}
\usepackage{array}
\usepackage{hyperref}
\hypersetup{colorlinks=false, pdfborder={0 0 0}}
\usepackage{fancyhdr}

\pagestyle{fancy}
\fancyhf{}
\fancyhead[L]{\small TrackFleet --- Papéis e Matriz RACI}
\fancyhead[R]{\small\thepage}
\renewcommand{\headrulewidth}{0.4pt}

\newcolumntype{L}{>{\raggedright\arraybackslash}X}
\newcolumntype{C}{>{\centering\arraybackslash}X}
\newcolumntype{R}{>{\raggedleft\arraybackslash}X}
\setlist{topsep=3pt}

% Cabeçalho padrão do projeto. No documento consolidado (trabalho_pratico_01/)
% estes três comandos são redefinidos e este mesmo arquivo vira parte do capítulo do domínio.
\newcommand{\cabecalho}[3]{%
  \begin{center}
  {\Large\bfseries DOCUMENTO DE #1 DO PROJETO}\\[2pt]
  {\large\bfseries TrackFleet --- Gestão de Rastreador de Automóveis}\\[4pt]
  Disciplina de Gerenciamento de Projetos $\cdot$ Framework PMBOK\textregistered\ 8\textsuperscript{a} Edição $\cdot$ Domínio: #2\\[2pt]
  Integrantes: Enzo Allebrand e Kerlon Ribeiro (dois GPs) $\cdot$ 4\textsuperscript{o} semestre $\cdot$ Data: #3
  \end{center}
  \vspace{2mm}\hrule\vspace{4mm}}
\newcommand{\parte}[1]{\noindent\textbf{\large #1}\par\vspace{1mm}}
\newcommand{\separador}{\par\vspace{2mm}\hrule\vspace{4mm}}

\begin{document}

\cabecalho{GOVERNANÇA}{Governança}{13/08}

\parte{Papéis e Matriz RACI}
\noindent Quem faz o trabalho e quem responde por cada decisão-chave do projeto, para que nenhuma decisão fique sem dono.

\section{Papéis}
\begin{itemize}[itemsep=3pt]
  \item \textbf{GP --- Enzo Allebrand:} conduz o projeto em co-gestão e responde pelas decisões técnicas de backend, motor de alertas e infraestrutura.
  \item \textbf{GP --- Kerlon Ribeiro:} conduz o projeto em co-gestão e responde pelas decisões técnicas de frontend, integração com o provedor de rastreamento, testes e LGPD.
  \item \textbf{Patrocinador:} Diretoria da transportadora parceira; aprova o termo, os marcos e as mudanças relevantes, e desempata impasses (papel fictício definido pela dupla).
\end{itemize}

\section{Matriz RACI}
R = faz $\cdot$ A = responde, só um por linha $\cdot$ C = consultado $\cdot$ I = informado $\cdot$ A/R = responde e também executa.

\begin{center}
\begin{tabularx}{\textwidth}{L c c c}
\toprule
\textbf{Decisão-chave} & \textbf{Patrocinador} & \textbf{GP (Enzo)} & \textbf{GP (Kerlon)} \\
\midrule
Aprovar a ideia / termo de abertura & A & R & R \\
Aprovar mudança de escopo, prazo ou objetivo & A & R & R \\
Decisões técnicas de backend, alertas e infraestrutura & I & A/R & C \\
Decisões técnicas de frontend, integração com o provedor, testes e LGPD & I & C & A/R \\
Homologar a entrega de cada etapa & A & R & R \\
\bottomrule
\end{tabularx}
\end{center}

\section{Observação}
Como os dois integrantes são GPs em co-gestão, a autoridade técnica do dia a dia foi dividida por frente: cada GP é o único responsável final (A) pelas decisões da sua frente e é consultado nas da frente do outro. As decisões que mudam o projeto (termo, mudanças relevantes e homologação) têm o patrocinador como único A. Assim, toda linha tem exatamente um responsável final, como exige a regra de ouro da RACI. A mesma divisão por frente aparece nos donos das atividades do Cronograma e na RACI detalhada do domínio de Recursos.

\end{document}
